\documentclass[letterpaper,11pt]{article}
\usepackage[T1]{fontenc}
\usepackage[utf8]{inputenc}
\usepackage[margin=1in]{geometry}
\usepackage[hidelinks]{hyperref}
\usepackage{xcolor}
\usepackage[default]{sourcesanspro}
\usepackage{parskip}
\usepackage[protrusion=true,expansion=false]{microtype}
%----------COLORS----------
\definecolor{ResumeNavy}{HTML}{18324A}
\definecolor{ResumeSlate}{HTML}{5B6573}
\definecolor{ResumeText}{HTML}{111827}
\color{ResumeText}
\setlength{\parskip}{8pt}
\pagestyle{empty}
\begin{document}
%----------HEADER (matches resume identity)----------
\begin{center}
  {\LARGE \bfseries \textcolor{ResumeNavy}{Deva Anand}} \\ \vspace{3pt}
  \small
  (413) 315-1715 \enspace\textbar\enspace
  \href{mailto:devaanand@umass.edu}{devaanand@umass.edu} \enspace\textbar\enspace
  \href{https://devaanand.com}{devaanand.com} \enspace\textbar\enspace
  \href{https://linkedin.com/in/devaaa/}{linkedin.com/in/devaaa} \enspace\textbar\enspace
  \href{https://github.com/Deva-1903}{github.com/Deva-1903}
\end{center}
\vspace{4pt}
{\color{ResumeNavy}\hrule height 1pt}
\vspace{12pt}

May 24, 2026

Snowflake University Recruiting \\
{\color{ResumeSlate}\textit{Re: Software Engineering Internship --- Fall 2026}}

\vspace{4pt}

Dear Snowflake Hiring Team,

I am a Master's student in Computer Science at UMass Amherst (graduating May 2027), applying for your Fall 2026 Software Engineering Internship. What draws me to this role specifically is how it is framed: delivering quality software in a large distributed system, with real emphasis on functional and performance testing, automation, and reliability. That is the kind of work I gravitate toward --- where correctness under load is the job, not overhead --- and it is exactly why a data platform that real customers depend on interests me far more than a greenfield side project.

Before grad school I spent two years owning a multi-tenant NHS clinical eTriage backend serving 8+ UK clients and 10,000+ monthly submissions. I treated integration tests, authentication hardening, release monitoring, and 20+ VAPT remediations as the core of the role, because a small mistake on sensitive clinical data genuinely cost something. More recently I optimized a Spark ETL pipeline over a 1.4-billion-row dataset, profiling stage-level timings in the Spark UI to localize shuffle and skew bottlenecks, cutting end-to-end runtime by roughly 25\%, and documenting a repeatable optimization playbook. Both experiences map closely to what your Quality \& Release Engineering and Database Query Engine teams care about: large-scale data, performance, and trustworthy software.

Your description of treating AI as a high-trust collaborator also resonates with how I already work. I use tools like Claude Code and Cursor as accelerators, but I keep system behavior grounded in tests and verifiable evidence --- my agentic web-discovery project, for instance, pairs LLM extraction with deterministic validation, cell-level provenance, and 150+ automated tests. I learn fastest when I am a little out of my depth and pick up unfamiliar tools quickly, so a fast-moving, experimental environment is where I do my best work. I would be glad to bring that mindset to a real system Snowflake customers rely on.

Thank you for your time and consideration.

\vspace{6pt}
Sincerely, \\
\textbf{Deva Anand}

\end{document}
